\section{Background and Related Work}

Six research lines touch parts of the policy pipeline, but they answer different questions.

\paragraph{Access-control analysis and mutation testing.}
Margrave computes semantic differences between two access-control policies and can explain requests whose decisions change \citep{margrave}. XACML mutation work defines policy fault operators, test adequacy, and equivalent-mutant handling \citep{xacml_mutation}. Cedar compares an authorization implementation with an independently specified model through formal and differential testing \citep{cedar_vgd}. These methods can analyze a policy language more deeply than \policystrata{} and may prove or systematically test decision equivalence inside that language. They do not by themselves connect a model-visible metric alias, semantic plan, SQL lowering, runtime database role, and result-release decision. \policystrata{} borrows mutation adequacy and differential-oracle ideas, but makes the unit of analysis a heterogeneous transition in the policy pipeline.

\paragraph{Translation validation and property-preserving compilation.}
\policystrata{}'s compiler contract is an instance of translation validation: rather than proving a compiler correct once, check each individual compilation against its input \citep{translation_validation}. Secure-compilation work sharpens what such a check should preserve, distinguishing full abstraction from the weaker robust preservation of specific property classes across a compilation boundary \citep{secure_compilation}. Authorization-preserving lowering is deliberately the weaker, property-class form: we do not ask a compiler to preserve all observational equivalences, only the declared obligations of principal, tenant, purpose, policy version, lineage, and metric semantics. What \policystrata{} adds is neither a new preservation theorem nor a proof technique, but the observation that a data agent runs several such lowerings between representations that are separately versioned, so a preservation failure must be attributed to one transition rather than to the pipeline as a whole.

\paragraph{SQL checking and generated databases.}
Distilled test-suite accuracy constructs databases that distinguish candidate SQL from a reference query \citep{test_suite_accuracy}; VeriEQL and SpotIt+ reason about bounded query equivalence under supported SQL semantics and constraints \citep{verieql,spotit}. Such tools give stronger SQL-level semantic evidence than finite execution alone. They do not decide whether a query was exposed to the right principal, retained tenant identity and policy version, used authorized lineage, or crossed an allowed release boundary. Conversely, \policystrata{} does not prove general SQL equivalence: its database states are finite witnesses, and a supported SQL verifier could replace that one checker.

\paragraph{Database and agent enforcement.}
Beacon checks consistency within DBMS access-control metadata and behavior \citep{beacon}, and row-level security is the containment mechanism our database contract observes \citep{postgresql_rls}. Agent enforcement systems place deterministic policy or privilege controls around tool execution \citep{pcas,progent}; Cordon adds a transactional boundary that stages irreversible effects until a composed execution flow validates \citep{cordon}; C-Trace compiles a subset of GDPR into predicates over execution traces and rejects non-compliant actions at runtime \citep{c_trace}; Data Flow Control pushes tuple-level release policy into the DBMS through provenance-aware query rewriting \citep{data_flow_control}; and role-conditioned refusal work evaluates model-level compliance \citep{role_conditioned_refusals}. These systems can block behavior at runtime. \policystrata{} instead tests whether separately maintained enforcement surfaces still agree with their assigned contracts, and it reports an upstream compiler violation even when RLS safely contains it. The relationship is complementary in both directions: an enforcement layer is a surface \policystrata{} can observe, and \policystrata{} is a way to detect that such a layer has drifted from the specification it was deployed to enforce.

\paragraph{Agent conformance testing and drift.}
The closest recent work shares our premise that agent safety degrades between components rather than inside one. Constraint drift argues that safety-critical constraints lose force as they pass through memory, delegation, communication, tool use, and audit, and that asserting a constraint once is not the same as maintaining it \citep{constraint_drift}. \policystrata{} is a narrow, executable instance of that position for a single-agent data pipeline: six named surfaces, declared obligations per transition, and a checker that reports where an obligation stopped holding. AgentRFC extracts normative clauses from agent protocol specifications into a typed IR, model-checks them, and replays counterexamples against live SDKs \citep{agentrfc}; its extract-check-replay shape is close to ours, but its subject is protocol conformance between agents rather than policy obligations between representations inside one agent. TDAD treats prompts as compiled artifacts and applies hidden test splits and semantic mutation testing to agent specifications \citep{tdad}, which is the same adequacy argument we make for a cross-layer operator registry. AgentRaft detects data over-exposure by building a cross-tool call graph, tracking taint at runtime, and adjudicating with a model committee \citep{agentraft}; it covers tool-to-tool flows that \policystrata{} does not model, while \policystrata{} stays deterministic, needs no model at scoring time, and localizes to a versioned transition rather than to a flow.

\paragraph{Generation and task benchmarks.}
Spider, BIRD, DAB, BEAVER, and DAComp measure SQL or data-task capability \citep{spider,bird,dab,beaver,dacomp}, while $\tau$-bench and ST-WebAgentBench measure whether agents follow domain rules and safety policies during interaction \citep{taubench,stwebagentbench}. PICARD and structured decoders reduce malformed outputs \citep{picard,outlines,xgrammar,jsonschemabench}. Correct, rule-following, or grammar-valid output can still select the wrong tenant, implement the wrong business metric, or violate release policy once it is lowered to SQL. \policystrata{} therefore starts at typed semantic IR; stochastic natural-language-to-plan reachability is reported separately and is not part of its deterministic score. A recent security study of data agents supplies the independent fault vocabulary we test our registry against \citep{data_agents_under_attack}; LASM supplies a second, structural taxonomy derived from 116 agent-security papers \citep{lasm}.

\begin{center}
\scriptsize
\setlength{\tabcolsep}{2pt}
\begin{tabular}{@{}L{0.24\columnwidth}L{0.31\columnwidth}L{0.35\columnwidth}@{}}
\toprule
System or method & Evidence it can provide & Boundary relative to \policystrata{} \\
\midrule
Margrave & semantic policy-version difference & one access-control language; no SQL or release path \\
XACML mutation & test adequacy against policy faults & mutates one policy representation \\
Cedar VGD & model/engine differential evidence & authorization decisions, not business semantics \\
Translation validation & per-compilation correctness & one lowering; no versioned surface attribution \\
Distilled SQL suites & distinguishing database states & query results, not authority or lineage \\
SQL verifiers & bounded equivalence evidence & supported SQL fragment only \\
Beacon & DBMS catalog/behavior inconsistency & database internals, not upstream agent surfaces \\
Runtime monitors & enforcement on observed actions & no first-transition diagnosis by itself \\
AgentRFC & protocol conformance via IR and replay & between agents, not between pipeline representations \\
AgentRaft & tool-flow over-exposure detection & model-adjudicated flows, not versioned lowerings \\
\policystrata{} & cross-representation witness and attribution & contract- and taxonomy-relative; no global proof \\
\bottomrule
\end{tabular}
\end{center}

The table is a capability comparison, not an empirical ranking: these systems were not run as drop-in competitors because they accept different languages and prove or test different properties. A fair composition would use Margrave or Cedar as the authorization oracle, VeriEQL as the compiler-semantic checker, and database/runtime monitors as observed surfaces inside \policystrata{}'s transition harness. The novelty boundary is therefore narrow. \policystrata{} is not a new policy logic, SQL verifier, or reference monitor. It composes evidence across versioned, non-equivalent representations and assigns a failed obligation to the first transition responsible for preserving it.
